\hypertarget{appendix-cac01ebeeb11}{}\label{referencestart:appendix-cac01ebeeb11}
\section{Axial depth and Euclidean range}\label{appendix-cac01ebeeb11-axial-depth-and-euclidean-range}

\textbf{Curriculum depth:} Developed · \textbf{First seen:} Phase 11 · \textbf{Latest development:} Phase 11

\subsection{Definition and mechanism}\label{appendix-cac01ebeeb11-definition-and-mechanism}

Axial depth measures distance along the camera's forward axis; Euclidean range measures distance from the camera centre along the viewing ray. With a right-handed view frame looking along \(-Z\), positive axial depth is \(d=-z\), while range is \(r=\sqrt{x^2+y^2+z^2}\). They agree only on the optical axis. For a unit ray direction \(q\) with \(q_z<0\), a range measurement converts to axial depth as \(d=-rq_z\). A range sensor and a perspective depth buffer may therefore encode different quantities. Applying a depth-buffer formula directly to a range image distorts reconstructed geometry away from the image centre. At view point \((3,4,-12)\), axial depth is \(12\) while range is \(13\), a visible distinction off-axis.

\subsection{Prerequisites}\label{appendix-cac01ebeeb11-prerequisites}

\begin{itemize}
\tightlist
\item
  \hyperref[appendix-1c07d64f5e70]{Camera view and clip-space pipeline}
\item
  \hyperref[appendix-f744e41b95b4]{Unit directions and safe normalization}
\end{itemize}

\subsection{Used by}\label{appendix-cac01ebeeb11-used-by}

\begin{itemize}
\tightlist
\item
  \hyperref[appendix-1bc40b08db4a]{Hardware depth-buffer linearization}
\item
  \hyperref[appendix-2cbf2968e866]{Position reconstruction from depth}
\item
  \hyperref[appendix-0976e86e5b9f]{Off-axis stereo projection and disparity}
\end{itemize}

\subsection{Full teaching explanation}\label{appendix-cac01ebeeb11-full-teaching-explanation}

\begin{itemize}
\tightlist
\item
  \href{../03-Geometry-Cameras-and-Image-Formation/main.pdf\#nameddest=02-foundations-ii-ch035-recover-geometry-from-depth-and-construct-screen-space-rays}{Recover geometry from depth and construct screen-space rays (Geometry, Cameras, and Image Formation, Geometry, Cameras, and Image Formation, Chapter 4)}
\end{itemize}

\subsection{Source coverage}\label{appendix-cac01ebeeb11-source-coverage}

Supplied lesson parts: 1078.

\subsection{Future keywords}\label{appendix-cac01ebeeb11-future-keywords}

range image, optical axis, axial distance
